% Part: proof-theory
% Chapter: natural-deduction
% Section: sequents

\documentclass[../../../include/open-logic-section]{subfiles}

\begin{document}

\olfileid[hi]{pt}{nat}{seq}

\olsection{अनुक्रमों सहित प्राकृतिक निगमन: \Log{N2c} और~\Log{N2i}}

\Log{G1c} या \Log{G1i} में !!^a{proof}~$\delta$ दिखाता है कि उसका
अंतिम-!!{formula}~$!A$ उसकी खुली मान्यताओं से निष्पन्न होता है। यदि $\Gamma$,
!!{formula}s~$!B$ का वह समुच्चय है जहाँ $!B^x$,~$\delta$ में खुली मान्यता
के रूप में आता है, तो इसे $\delta \Proves \Gamma \Sequent !A$ लिख सकते हैं।
इससे संकेत मिलता है कि हम अनुक्रमों के लिए भी प्राकृतिक निगमन प्रणाली
सूत्रबद्ध कर सकते हैं। ऐसी प्रणाली में प्रत्येक अनुक्रम यह जानकारी रखता है
कि दाईं ओर का !!{formula} किन खुली मान्यताओं पर निर्भर है, अर्थात् अनुक्रम पर
समाप्त होने वाले उप-!!{proof} में कौन-सी मान्यताएँ खुली हैं। यह अनुक्रमों
सहित प्राकृतिक निगमन, या ``अनुक्रम शैली में'' प्राकृतिक निगमन है, और हम
फलित प्रणालियों को \Log{N2c} तथा~\Log{N2i} कहते हैं।

\Log{N1} और \Log{N2} के बीच संगति को अधिक निकट बनाने के लिए बाईं ओर के
समुच्चयों~$\Gamma$ को चिह्नित !!{formula}s~$x:!B$ से बना लें। चूँकि \Log{N1}
के नियम केवल आधार-वाक्यों के !!{formula}s, अर्थात् निकटस्थ उप-!!{proof}s के
अंतिम-!!{formula}s, पर काम करते हैं, इसलिए \Log{N2} के नियम केवल अनुक्रम के
अनुवर्ती पर काम करते हैं। अतः अनुक्रमों के बाईं ओर के चिह्नित सूत्रों को हम
केवल प्रसंग कह सकते हैं।

मान्यताओं के स्थान पर \Log{N2} प्रमाणों के सबसे ऊपर के अनुक्रम इस रूप के
अभिगृहीत होते हैं
\[x:!A \Sequent !A.\]
नियम ठीक \Log{N1} के नियमों के अनुरूप हैं और \olref{tab:N2} में दिए गए हैं।

\subfile{rules-N2}

क्योंकि हमें (\Log{N1c} और \Log{N1i} की तरह) यह अनुमति देनी है कि नियम
\Intro{\lif}, \Intro{\lnot}, \Elim{\lor} और \Elim{\lexists} मान्यताओं को
वास्तव में मुक्त कर सकते हैं, पर ऐसा आवश्यक नहीं, इसलिए हम \Log{N2c}
और~\Log{N2i} में संबद्ध नियमों के सही अनुप्रयोग तब भी स्वीकार करते हैं जब
संबद्ध प्रसंग !!{formula}s $!A$, $!B$ और $!A(a)$ संबद्ध आधार-वाक्य के प्रसंग
में उपस्थित न हों।

\begin{prop}
यदि $\delta$,~$!A$ का \Log{N1c} या \Log{N1i} में कोई !!a{proof} है, तो
~\Log{N2c} (\Log{N2i}) में $\Gamma \Sequent !A$ का कोई !!a{proof}~$\delta'$
है, जहाँ $\Gamma$ उन सभी~$x:!B$ का समुच्चय है जिनके लिए $!B^x$,~$\delta$ की
खुली मान्यता है।
\end{prop}

\begin{proof}
  $n=\pheight{\delta}$ पर आगमन से। यदि $n=0$, तो $\delta$ अकेली खुली
  मान्यता~$!A^x$ है। तब $\Gamma = \{x:!A\}$ और $\Gamma \Sequent !A$,
  \Log{N2c} या~\Log{N2i} का अभिगृहीत है।

  यदि $n>0$, तो हम~$\delta$ के अंतिम अनुमान के अनुसार स्थितियों में भेद करते
  हैं। हम केवल एक स्थिति की चर्चा करते हैं; शेष अभ्यास के रूप में छोड़ी जाती
  हैं।

  मान लें अंतिम अनुमान $\Intro{\lif}$ है। तब $\delta$ का अंत है
  \[
  \AxiomC{$\Discharge{!B}{x}$}
  \RightLabel{$\delta_1$}
  \DeduceC{$!C$}
  \DischargeRule{\Intro{\lif}}{x}
  \UnaryInfC{$!B \lif !C$}
  \DisplayProof
  \]
  आगमन परिकल्पना से~\Log{N2c} (\Log{N2i}) में $\Gamma' \Sequent !C$ का कोई
  !!a{proof}~$\delta_1'$ है, जहाँ $\Gamma'$ वे चिह्नित मान्यताएँ हैं जो
  ~$\delta_1$ में खुली हैं। यदि $!B^x$, $\delta_1$ की खुली मान्यता है, तो उसे
  \Intro{\lif} अनुमान पर मुक्त किया जाता है और~$\delta$ की खुली मान्यताएँ
  $\Gamma = \Gamma' \setminus \{x:!B\}$ हैं। दूसरे शब्दों में $\delta_1'$,
  ~ $x:!B, \Gamma \Sequent !C$ का !!a{proof} है और हम $\delta'$ को यह ले
  सकते हैं
  \[
  \AxiomC{}
  \RightLabel{$\delta_1'$}
  \Deduce$x:B, \Gamma \fCenter !C$
  \RightLabel{\Intro{\lif}}
  \UnaryInf$\Gamma \fCenter !B \lif !C$
  \DisplayProof
  \]
  यदि $!B^x$,~$\delta_1$ की खुली मान्यता नहीं है, तो \Intro{\lif} कोई
  मान्यता मुक्त नहीं करता और $\delta$ तथा $\delta_1$ की खुली मान्यताएँ समान
  हैं। चूँकि~\Log{N2c} और~\Log{N2i} में \Intro{\lif} के आधार-वाक्य में
  प्रसंग सूत्र $x:!B$ का उपस्थित होना आवश्यक नहीं, हम $\delta'$ को यह ले
  सकते हैं
  \[
  \AxiomC{}
  \RightLabel{$\delta_1'$}
  \Deduce$\Gamma \fCenter !C$
  \RightLabel{\Intro{\lif}}
  \UnaryInf$\Gamma \fCenter !B \lif !C$
  \DisplayProof
  \]

  वे स्थितियाँ भी इसी प्रकार हैं जहाँ $\delta$ किसी अन्य नियम पर समाप्त होता
  है।
\end{proof}

\begin{prop}
यदि $\delta$, $\Gamma \Sequent !A$ का \Log{N2c} या \Log{N2i} में कोई
!!a{proof} है, तो \Log{N1c} (\Log{N1i}) में कोई !!a{proof}~$\delta'$ है
जिसका अंतिम-!!{formula}~$!A$ और प्रत्येक $x:!B \in \Gamma$ के लिए खुली
मान्यता $!B^x$ है।
\end{prop}

\begin{proof}
  हम फिर~$\delta$ की !!{height}~$n$ पर आगमन से आगे बढ़ते हैं। यदि $n=0$,
  तो $\delta$ अभिगृहीत $x:!A \Sequent !A$ है। !!{proof}~$\delta'$ मान्यता
  ~$!A^x$ है।

  यदि $n>0$, तो हम~$\delta$ के अंतिम अनुमान के अनुसार स्थितियों में भेद करते
  हैं। उदाहरणतः यदि वह अनुमान~$\Intro{\lif}$ है, तो $\delta$ का अंत है
  \[
  \bottomAlignProof
  \AxiomC{}
  \RightLabel{$\delta_1$}
  \Deduce$x:!B, \Gamma \fCenter !C$
  \RightLabel{\Intro{\lif}}
  \UnaryInf$\Gamma \fCenter !B \lif !C$
  \DisplayProof
  \quad\text{ या }\quad
  \bottomAlignProof
  \AxiomC{}
  \RightLabel{$\delta_1$}
  \Deduce$\Gamma \fCenter !C$
  \RightLabel{\Intro{\lif}}
  \UnaryInf$\Gamma \fCenter !B \lif !C$
  \DisplayProof
  \]
  आगमन परिकल्पना से~\Log{N1c} (\Log{N1i}) में~$!C$ का कोई !!a{proof} है।
  पहली स्थिति में $!B^x$,~$\delta'$ की खुली मान्यता है; दूसरी में नहीं। हम
  !!{proof}~$\delta'$ को यह ले सकते हैं
  \[
  \bottomAlignProof
  \AxiomC{$\Discharge{!B}{x}$}
  \RightLabel{$\delta_1$}
  \DeduceC{$!C$}
  \DischargeRule{\Intro{\lif}}{x}
  \UnaryInfC{$!B \lif !C$}
  \DisplayProof
  \quad\text{ या }\quad
  \bottomAlignProof
  \AxiomC{}
  \RightLabel{$\delta_1$}
  \DeduceC{$!C$}
  \RightLabel{\Intro{\lif}}
  \UnaryInfC{$!B \lif !C$}
  \DisplayProof
  \]
  क्रमशः।
\end{proof}

\end{document}
